% =====================================================================
\chapter{Full Result Tables}
\label{app:tables}
% =====================================================================

This appendix gives the complete per-benchmark measurements behind the summary figures and
tables of Chapter~\ref{ch:results}. All numbers are AUROC, read directly from the saved
measurement files; $0.5$ is chance.

Table~\ref{tab:app-gptj-full} is the full matrix for the six-billion-parameter model that
anchors the study: every one of the sixteen configurations and six baselines, scored on the
same rows of all ten benchmarks. It is the table from which the means in
Section~\ref{sec:gptj} are computed, and it makes the per-benchmark spread --- in particular
the dominance of the HaluEval-QA column --- explicit.

\begin{table}[h]
\centering
\caption{Full per-benchmark transfer AUROC on the six-billion-parameter model (GPT-J): sixteen feature configurations and six baselines, scored on identical rows of all ten benchmarks. \textbf{Bold} marks the best value in each column. Source: the consolidated configuration and baseline measurement files.}
\label{tab:app-gptj-full}
\footnotesize
\setlength{\tabcolsep}{3.2pt}
\begin{tabular}{lcccccccccc|c}
\toprule
Method & TruthQA & TrivQA & CoQA & TydiQA & HE-QA & HE-Sm & HE-Dl & NQ & HotQA & PopQA & Avg \\
\midrule
C1 canonical & 0.503 & 0.476 & 0.580 & 0.502 & 0.124 & 0.398 & 0.330 & 0.562 & 0.569 & 0.541 & 0.459 \\
C2 {+}drift & 0.492 & 0.459 & 0.566 & 0.492 & 0.127 & 0.337 & 0.339 & 0.551 & 0.558 & 0.553 & 0.447 \\
C3 {+}variance & 0.478 & 0.453 & 0.564 & 0.497 & 0.109 & 0.381 & 0.349 & 0.535 & 0.554 & 0.534 & 0.445 \\
C4 {+}entropy & 0.519 & 0.410 & 0.541 & 0.478 & 0.260 & 0.389 & 0.389 & 0.513 & 0.524 & \textbf{0.692} & 0.472 \\
C5 {+}trio & 0.526 & 0.503 & 0.570 & 0.487 & 0.272 & 0.371 & 0.330 & 0.403 & 0.557 & 0.620 & 0.464 \\
C6 {+}lookback & 0.498 & 0.483 & 0.591 & 0.505 & 0.068 & 0.342 & 0.321 & 0.544 & 0.550 & 0.548 & 0.445 \\
C7 {+}logit-lens & 0.472 & 0.482 & 0.573 & 0.502 & 0.079 & 0.383 & 0.344 & 0.535 & 0.546 & 0.545 & 0.446 \\
C8 {+}margin & 0.522 & 0.478 & 0.525 & 0.536 & 0.434 & 0.374 & 0.396 & 0.558 & 0.517 & 0.417 & 0.476 \\
C9 {+}l/l/m & 0.504 & 0.511 & 0.563 & 0.490 & 0.126 & 0.348 & 0.329 & 0.455 & 0.563 & 0.569 & 0.446 \\
C10 {+}all lit. & 0.494 & 0.540 & 0.454 & 0.501 & 0.844 & 0.597 & \textbf{0.668} & 0.513 & 0.446 & 0.426 & 0.548 \\
C11 {+}trio{+}llm & 0.484 & 0.518 & 0.546 & 0.520 & 0.635 & 0.391 & 0.464 & 0.539 & 0.603 & 0.553 & 0.525 \\
C12 {+}trio{+}lit. & 0.555 & 0.534 & 0.444 & 0.465 & 0.408 & 0.462 & 0.522 & 0.484 & 0.453 & 0.355 & 0.468 \\
C13 {+}urp/curv & 0.494 & 0.507 & 0.576 & 0.507 & 0.065 & 0.315 & 0.311 & 0.572 & 0.567 & 0.596 & 0.451 \\
C14 {+}conf-traj & 0.481 & 0.497 & 0.421 & 0.493 & 0.944 & 0.564 & 0.665 & 0.516 & 0.427 & 0.451 & 0.546 \\
C15 {+}urp/cv/cf & 0.504 & 0.481 & 0.440 & 0.461 & 0.351 & 0.527 & 0.406 & 0.604 & 0.449 & 0.483 & 0.471 \\
C16 {+}trio{+}new & 0.465 & 0.471 & 0.475 & 0.475 & 0.620 & 0.580 & 0.585 & 0.522 & 0.466 & 0.424 & 0.508 \\
\midrule
\multicolumn{12}{l}{\emph{Baselines}}\\
SAPLMA & 0.470 & 0.486 & 0.515 & 0.534 & \textbf{0.985} & \textbf{0.611} & 0.649 & 0.471 & 0.555 & 0.341 & \textbf{0.562} \\
Perplexity & \textbf{0.570} & 0.580 & \textbf{0.609} & \textbf{0.600} & 0.434 & 0.397 & 0.426 & 0.589 & 0.597 & 0.429 & 0.523 \\
HalluShift & 0.466 & 0.512 & 0.558 & 0.547 & 0.446 & 0.474 & 0.469 & 0.612 & 0.578 & 0.430 & 0.509 \\
HaloScope & 0.523 & 0.472 & 0.512 & 0.545 & 0.455 & 0.531 & 0.535 & \textbf{0.627} & 0.534 & 0.352 & 0.508 \\
MIND & 0.522 & 0.462 & 0.542 & 0.458 & 0.520 & 0.384 & 0.505 & 0.532 & 0.493 & 0.453 & 0.487 \\
EigenScore & 0.523 & \textbf{0.641} & 0.400 & 0.402 & 0.501 & 0.500 & 0.500 & 0.395 & \textbf{0.656} & 0.347 & 0.487 \\
\bottomrule
\end{tabular}
\end{table}

Table~\ref{tab:app-opt-diag} is the file-backed diagnostic run for the second
six-billion-class model. It covers twelve configurations on the original seven-benchmark
suite at a small number of rows per benchmark; as noted in Chapter~\ref{ch:results}, these
numbers are diagnostic rather than scientific, and the extended ten-benchmark run with
baselines is awaiting archival.

\begin{table}[h]
\centering
\caption{Diagnostic configuration run on the second six-billion-class model (OPT-6.7B), the part backed by a measurement file in the repository: twelve configurations on the original seven-benchmark suite at a small number of rows per benchmark. At this sample size the numbers are diagnostic, not scientific. AUROC; $0.5$ is chance.}
\label{tab:app-opt-diag}
\footnotesize
\setlength{\tabcolsep}{4pt}
\begin{tabular}{lccccccc|c}
\toprule
Cfg. & TruthQA & TrivQA & CoQA & TydiQA & HE-QA & HE-Sm & HE-Dl & Avg \\
\midrule
C1 canonical & 0.328 & 0.535 & 0.571 & 0.503 & 0.054 & 0.423 & 0.345 & 0.394 \\
C2 {+}drift & 0.338 & 0.567 & 0.564 & 0.523 & 0.039 & 0.316 & 0.313 & 0.380 \\
C3 {+}variance & 0.308 & 0.561 & 0.550 & 0.537 & 0.025 & 0.321 & 0.310 & 0.373 \\
C4 {+}entropy & 0.303 & 0.539 & 0.556 & 0.542 & 0.126 & 0.342 & 0.336 & 0.392 \\
C5 {+}trio & 0.276 & 0.485 & 0.502 & 0.492 & 0.228 & 0.380 & 0.342 & 0.386 \\
C6 {+}lookback & 0.286 & 0.546 & 0.558 & 0.506 & 0.032 & 0.361 & 0.316 & 0.372 \\
C7 {+}logit-lens & 0.313 & 0.527 & 0.542 & 0.514 & 0.026 & 0.336 & 0.309 & 0.367 \\
C8 {+}margin & 0.308 & 0.528 & 0.559 & 0.516 & 0.037 & 0.335 & 0.313 & 0.371 \\
C9 {+}l/l/m & 0.303 & 0.519 & 0.551 & 0.504 & 0.055 & 0.334 & 0.349 & 0.374 \\
C10 {+}all lit. & 0.308 & 0.626 & 0.562 & 0.567 & 0.246 & 0.395 & 0.353 & 0.437 \\
C11 {+}trio{+}llm & 0.328 & 0.553 & 0.551 & 0.538 & 0.049 & 0.286 & 0.302 & 0.373 \\
C12 {+}trio{+}lit. & 0.293 & 0.528 & 0.554 & 0.498 & 0.020 & 0.352 & 0.318 & 0.366 \\
\bottomrule
\end{tabular}
\end{table}

% =====================================================================
\chapter{Hyperparameters and Implementation Details}
\label{app:hyper}
% =====================================================================

For reproducibility, this appendix records the settings used throughout. All runs use a
fixed random seed of $42$.

\section*{Data generation}
\begin{itemize}[leftmargin=1.6em,itemsep=0.2em]
\item Source text: encyclopedia articles; the first sentences are retained and a named
entity occurring after the first sentence is selected.
\item Faithful/hallucinated decision: the original entity's first token must appear among
the model's top four predicted tokens to be labelled faithful.
\item Generation cap during labelling: a short continuation of up to $48$ new tokens.
\item Scale: one thousand examples per class (two thousand total) for the large models;
smaller counts for the iteration and smoke models, as recorded in Chapter~\ref{ch:results}.
\item Numeric format: bfloat16, with no further quantisation.
\end{itemize}

\section*{Classifier}
\begin{itemize}[leftmargin=1.6em,itemsep=0.2em]
\item Architecture: a five-layer perceptron with progressively narrowing hidden widths,
rectified-linear activations, and a two-way softmax output.
\item Loss and optimiser: cross-entropy; Adam with learning rate $5\times10^{-4}$ and weight
decay $10^{-5}$.
\item Batch size $32$; ten training epochs.
\item Features standardised by a scaler fitted on the training split only.
\item Split: $80\%$ train / $20\%$ test, identical across all configurations.
\end{itemize}

\section*{Evaluation}
\begin{itemize}[leftmargin=1.6em,itemsep=0.2em]
\item Each method --- configuration or baseline --- is scored on the same deterministic
first-$N$ slice of each benchmark, with $N$ up to $350$ rows per benchmark in the unified
runs.
\item Metrics: AUROC (primary), accuracy, precision, recall, F1; calibration (Brier score,
expected calibration error) where the measurement file records it.
\item In-distribution and transfer numbers are never averaged together.
\end{itemize}

\section*{Baseline settings (six-billion model)}
\begin{itemize}[leftmargin=1.6em,itemsep=0.2em]
\item Supervised mid-layer probe: hidden state read at an intermediate layer (layer $17$ of
$28$).
\item Logit-lens divergence feature: early/late comparison taken at layer $14$.
\item Spectral baseline: best layer and threshold selected on a validation split.
\item Sampling-based score: a small number of stochastic samples per query, with a ridge
term for numerical stability.
\end{itemize}

\section*{Hardware}
\begin{itemize}[leftmargin=1.6em,itemsep=0.2em]
\item Primary: a free cloud notebook with two sixteen-gigabyte T4 GPUs (thirty-two gigabytes
combined), nine-hour session cap, thirty GPU-hours per week.
\item Fallback: a free single-T4 notebook for the smaller models.
\item The seven-billion models require both GPUs; a single sixteen-gigabyte card cannot hold
a seven-billion model in bfloat16 together with generation memory.
\end{itemize}

% =====================================================================
\chapter{Reproducibility Notes}
\label{app:repro}
% =====================================================================

The study is designed to be reproducible on a free account. Three points are worth
emphasising. First, the training data is generated separately for each model, because probes
do not transfer between models; a detector is therefore always trained on data from the same
model it monitors. Second, every configuration and baseline within a model shares the same
generated data, the same train/test split and seed, and the same evaluation rows, so that
differences in measured performance are attributable to the features rather than to incidental
differences in setup. Third, the evaluation benchmarks are downloaded once and cached so that
runs are not dependent on network conditions, and the per-benchmark row count is fixed so that
cost stays within the free-tier budget. Where a model's measurement file is not yet archived,
its numbers are marked as pending in Chapter~\ref{ch:results} rather than reported as results.
